\uuid{cVid}
\titre{Espaces vectoriels}
\auteur{hanani ; mzari}
\organisation{exo7}
\licence{CC BY-NC-SA 4.0}
\datecreate{2020-01-20}
\niveau{L1}
\module{Algèbre}
\chapitre{Espace vectoriel}
\sousChapitre{Définition, sous-espace}
\difficulte{2}
\theme{Espaces vectoriels}
\exo7qcmid{lille-2.espaces-vectoriels.012}

\contenu{
\begin{question}
\qclassification{106.01}
\qtags{motcle=Sous-espaces vectoriels de $\Rr^n$}
Quelles sont les assertions vraies ?
\begin{answers}
\bad{Les sous-espaces vectoriels de $\Rr^2$ sont les droites vectorielles.}
\good{Les sous-espaces vectoriels non nuls de $\Rr^2$ sont les droites vectorielles et $\Rr^2$.}
\bad{Les sous-espaces vectoriels de $\Rr^3$ sont les plans vectoriels.}
\good{Les sous-espaces vectoriels non nuls de $\Rr^3$  qui sont strictement inclus dans $\Rr^3$ sont les droites vectorielles et les plans vectoriels.}
\end{answers}
\begin{explanations} Les sous-espaces vectoriels de $\Rr^2$ sont : $\{(0,0)\}$, les droites vectorielles et $\Rr^2$. Les sous-espaces vectoriels de $\Rr^3$ sont : $\{(0,0,0)\}$, les droites vectorielles, les plans vectoriels et $\Rr^3$.
\end{explanations}
\end{question}
}
